\section{鲁、郑、宋：中原小国怎样活在大国之间}

洛邑以东、以南的道路把鲁、郑、宋连入同一片诸侯世界，却没有给它们同一种命运。鲁在曲阜，国力未必压过邻人，留下的年次却成为后人辨认春秋的经线；郑贴着王畿，既替周室办事，也最先感到王命不能再直接调动地方；宋夹在晋、楚之间，兵车往来时常是受压的一方，僵局形成时又可能成为双方都肯入席的地方。小国的作为不在于脱离大国，而在于看准道路、礼仪和资源暂时留下的缝隙。

\begin{event}{公元前722年}{鲁隐公元年与《春秋》年表开始}{可靠纪年}{鲁}\eventanchor{SA-0722-CHUNQIU}{春秋·鲁隐公元年}
曲阜的史官以“元年春王正月”落下第一笔。往后的字句短得近乎冷峻：会盟、来朝、伐国、葬君、日食，都被安进鲁君在位的年次。它用的是周王正月，站的却是鲁国朝廷的门槛；列国由此有了一条可以相互校验的时间线，后世也因此容易把鲁人看见的事误作整个天下唯一的视野。

这种可见性本身是一种政治资源。东迁后的王室尚能提供月令、爵名和礼仪的共同语言，却已不能替所有诸侯裁断争端。鲁国没有因保存年表便取得霸权；它所保住的，是把各国行动写成先后相承之事的权力。经文中看似平静的称谓、动词和取舍，正是后来的传家、经师和史家不断进入的门。

\term{史料边界}：隐公元年的经文是可靠的年代起点。《左传》《公羊传》《谷梁传》分别把这部短书展开为叙事、义例和经学解释；它们不能彼此自动证实。孔子删修《春秋》是影响深远的传统说法，但现存经文的档案基础、编纂层次和传抄过程仍是研究问题，不能把每一字都当作一位作者留下的现场笔录。
\end{event}

鲁史开篇的同一年，郑国已经显出另一种小国难题。郑庄公即位后，母族、弟弟共叔段和封邑之间的冲突终于公开。经文以“\eventanchor{SA-0722-ZHENG-DUAN}{春秋·郑伯克段于鄢}”记其事，留下的是宗君击败弟弟的结果；《左传》才铺开武姜偏爱、请封、筑城、庄公等待其变，以及“克段于鄢”的道德张力。可以确知的是，继承竞争和封邑所附的人力、城邑使郑的近畿权力不能自然集中；至于母子之间每一次言语、庄公是否预作圈套，属于后出叙事最浓的部分。

庄公取胜并没有使郑成为脱离周室的独立堡垒，反而使他更能以卿士身份同王室周旋。郑武公、庄公父子长期居于周朝政务的要位，郑既可调用本国的邑兵与粮食，也要为王室的衰弱承担距离最近的压力。平王晚年改倚虢公，周、郑互换人质，郑人又取温地之麦、成周之禾；《左传》把这些事连成一条逼向战争的线索。它提示冲突不只是抽象的“尊王”与“叛王”，还牵涉可耕之地、供军之粮和卿士职位；不过，具体交涉的先后和言辞仍须保留传文的叙事层次。

\begin{event}{公元前707年}{郑庄公与王师战于繻葛}{可靠纪年}{繻葛}\eventref{SA-0707-XUGE}{春秋·繻葛之战}
桓王终于率诸侯之师伐郑。繻葛的地望至今有不同考订，战场却足够靠近王畿，让这次交锋格外刺眼：王室以旧有的征发名义出兵，郑则以一国军政资源迎战。《春秋》只写“王师败绩于繻葛”，已把胜负钉在年表上；《左传》所写祝聃射中王肩、庄公战后处置，则把兵事编进礼与名分的戏剧性叙述，不能逐字当作战地实录。

郑军的胜利没有废去周王的名分。它显示的更为具体：当王室的兵源、财赋和指挥无法压倒近畿诸侯时，掌握城邑和军队的卿士可以拒命而仍保有位置。直接后果是王室军事威信受损；中期影响则是诸侯越来越依靠会盟、暂时从属和武力处理争端。郑的强势也有边界，它来自紧邻洛邑的特殊条件，不能无限外推为一种稳定霸权。
\end{event}

\begin{sourcenote}[郑庄公的家事与王事]
《春秋·隐公元年》与《春秋·桓公五年》分别固定了克段、王师败绩的年代骨架；《左传·隐公元年、三年及桓公五年》把封邑、质子、麦禾和箭伤编成连续故事，《史记·郑世家》又以汉人的叙述方式承接了这条线。互读可见郑的内政集中和周郑摩擦彼此相关，却不能把传文中完整的心理、对白和战术细节当作同时代材料已经独立证明的事实。
\end{sourcenote}

鲁也不是只在竹简上观看别国。前684年，齐军南来，鲁军在长勺迎战。\eventanchor{SA-0684-CHANGSHAO}{春秋·长勺之战}的经文确认齐伐鲁和鲁师获胜，说明曲阜仍能在边境保住自身；《左传》以曹刿论战的节奏解释“齐师败绩”，遂使鼓进鼓止成为后世最熟悉的春秋战场画面。那些言辞保存的是后人理解车战、军心与主将判断的方式，不等于我们能够据此复原每一步战术。鲁赢得的是一段边境喘息，而不是改写齐国动员能力的资本；齐的资源和联盟随后仍将转向更广阔的北方与中原。

宋的困局比鲁更裸露。齐桓公死后，宋襄公试图接过会盟的声望；前638年，他在泓水与楚军相遇。\eventanchor{SA-0638-HONG}{春秋·泓水之战}以宋师败绩收束，经文所能牢固确认的是战事、交战者和结果。后出的《左传》写宋襄公不肯趁楚军渡河、未成列时攻击，把一次败战塑成守礼而失机的寓言；《国语·宋语》与《史记·宋微子世家》又各自延续宋襄公的形象。与其把这个故事简单读作“仁义误国”，不如先看见宋所面对的资源差距：它想以礼仪号召诸侯，楚却能以北进的军政力量争夺同一批关系。泓水之后，宋的霸权尝试受挫，楚在中原的存在更难忽略，并接入了\eventref{SA-0632-CHENGPU}{春秋·城濮之战}前后的晋楚竞争。

宋人并未因此只能挨打。晋、楚连年争胜，宋恰在军队和使者必须经过的中间地带；降低战争风险首先是它自己的生存需要。前579年，宋大夫华元推动晋楚在宋结盟。\eventanchor{SA-0579-HUAYUAN}{春秋·宋华元调停}所见的，是一次并未终结竞争的停歇：两强都还保有兵力和附属国，也都暂时不愿继续承受征发的代价。《春秋·成公十二年》提供会盟的年次，《左传》和《国语·宋语》所存谈判言辞则带有很强的外交修辞。华元未能创造一个凌驾诸侯的仲裁者，却让“由宋居间、以会盟减损战事”成为可再使用的办法。

\begin{event}{公元前546年}{宋向戌调停晋楚于弭兵之会}{可靠纪年}{宋}\eventref{SA-0546-MIBING}{春秋·弭兵之会}
经过\eventref{SA-0575-YANLING}{春秋·鄢陵之战}，晋、楚都还能够动员，谁也难以用一场战役迫使对方退出中原。宋的执政大夫向戌在这道僵局中来往于两边；《左传·襄公二十七年》用“将平晋、楚”概括他的目标。调停者不是站在冲突之外的裁判，而是必须让两方相信：把会面放在宋，所得会比继续出兵更多。

会盟在宋举行，所达成的并不是现代意义的和平条约。晋、楚各自的从属网络被暂时并列承认，诸侯因而要承受双重而紧张的关系；齐、秦等国也没有因此纳入一个统一的指挥体系。宋获得了罕有的外交能见度，同时要承担接待、传递意向和维持中立的风险。它的力量不在替晋楚下令，而在于把双方不愿放弃的礼仪空间变成一张可以谈判的席位。

弭兵减少了晋楚直接决战，不能消除兼并，更没有设立一支能执行盟约的常备力量。外部冲突暂缓后，诸国卿族和执政集团反而有更多时间经营封邑、家臣和军队。它是消耗之后的均势安排，不是春秋秩序的终局；华元的先例和向戌的成功，恰好说明宋的主动何其依赖晋楚谁也不能独胜的时刻。
\end{event}

\begin{causalchain}[从王师败绩到宋地会盟]
郑在\eventref{SA-0707-XUGE}{春秋·繻葛之战}所揭示的，是王室已难独断近畿；到\eventref{SA-0546-MIBING}{春秋·弭兵之会}，问题已不再是恢复王命，而是怎样让并立的晋楚减少相互耗损。这不是从一场败仗直线推出一场会盟，而是旧有裁断机制逐渐失效后，各地分别以武力、盟约和调停填补空缺的过程。
\end{causalchain}

郑在这种两强夹缝中走得更久，代价是把许多政治争执带进城内。前543年前后，子产开始主掌郑国政务。\eventanchor{SA-0543-ZICHAN-POWER}{春秋·子产执政}并非一纸新制的突然降临：郑要在晋楚之间选择进退，又要处理卿族、城邑和诉讼，稳定可预期的行政比一时扩张更迫切。《左传·襄公三十年》《国语·郑语》与《史记·郑世家》所保存的子产言行，多已成为后世的政治记忆；可以用来观察小国治理的问题，不能把其中每段议论视为当日的速记。

前536年，郑人铸刑书，\eventanchor{SA-0536-ZHENG-XINGSHU}{春秋·子产铸刑书}把原先更多依赖贵族议礼的裁断规则以器物公布。《左传·昭公六年》还保存叔向致书批评的说法，使后世常把它读成礼与法的一次正面冲突。实际公布的条文、适用范围，以及书信的真伪和成文层次，学界仍有讨论；但这一举动至少表明，法律知识不再完全锁在少数人对礼例的掌握之中。它提高了裁断的可预期性，也可能使原有贵族对解释规则的垄断受到挑战。

前522年子产去世，\eventanchor{SA-0522-ZICHAN-DEATH}{春秋·子产去世}后的郑仍须在晋楚之间维持位置。传世材料中民众哀悼、后继者称誉子产的篇幅很动人，却是《左传》《国语》以政治理想回望一位执政者的方式。较能把握的是他的长期经营没有替郑消去外压，只是使这座近畿国家在资源有限时拥有了较稳的处置手段；后人所称道的子产，也由此成为评价小国政治的尺度，而不是一段可以脱离条件复制的答案。

鲁的“记录者”位置同样没有保护鲁君。三桓长期掌握邑、兵和朝政，鲁昭公试图反制而失败；前517年，\eventanchor{SA-0517-LU-ZHAO-EXILE}{春秋·鲁昭公出奔}的经文记下君主奔齐，鲁国公室的空心化遂无法再用礼仪遮蔽。《左传》关于攻郓等战事的展开，能让我们看见公室与季氏集团如何在城邑间争夺，细部仍属叙事材料，需要和经文的结果分开读。国君出奔并非鲁政一夕崩解，而是多年资源由公室流向卿族的公开显影。

五年后，季氏的家臣阳虎也在斗争中失势出奔。经文记下\eventanchor{SA-0502-YANGHU}{春秋·阳虎出奔}，让权力再向下移的一刻浮出纸面：公室受制于三桓，三桓又可能受制于掌兵掌邑的家臣。阳虎的动机及其同孔子的后出关联，不能以传说补成确定事实；但他的败走并未把权力送回鲁君，反而显示封邑、家臣和武装一旦形成网络，谁控制它们就可以暂时越过名分行事。

\begin{event}{公元前481年}{获麟与齐国陈氏的权力跃升}{可靠纪年}{鲁与齐}\eventanchor{SA-0481-LIN}{春秋·获麟与陈氏政变}
鲁哀公十四年，鲁国经文写下“西狩获麟”。四字像一扇忽然合上的门，后世遂常以\eventref{SA-0481-END}{春秋·记事终点}标识《春秋》的终篇。它首先是文本传承的终点，不是列国政治在这一年忽然换了一副面貌。

同年，齐国陈恒杀齐简公壬于舒州，改立其弟骜。经文称“陈恒”，后世通称“田恒”，是同一支陈、田氏族的不同书写；这次行动使田氏在齐国内部的控制越过又一道门槛。它仍不是\eventref{WQ-0386-TIANQI}{战国·田氏代齐}的提前完成。若把前481年直接叫作“田氏代齐”，便会抹去贵族家族长期积累封邑、家臣和军政资源的过程，也会把政变后的实际控制误作法统上的改朝换代。

获麟与陈恒弑君同年出现，并不意味着前者造成后者。鲁的记录停止，齐的卿族上升，却让人同时看见春秋末叶两种转折：一边是旧编年传统给出的收束，一边是国君与世卿之间的权力天平继续倾斜。此前弭兵保留的均势并未冻结各国内部的竞争；齐国的变化，正是这种长期重组在一国之内最显眼的一幕。
\end{event}

\begin{textwindow}[《春秋》的最后一笔]
“西狩获麟。”经文只有四字。《公羊传》把它解释为孔子“吾道穷矣”的象征；《左传》与《史记》的齐事则把注意力转向陈恒弑君及其后的权力安排。原文无需逐字翻译；需要分开的是经文的停止、经学传统赋予它的象征，以及齐国内部实际而漫长的家族权力转移。
\end{textwindow}

\begin{turningpoint}[制度得失：小国的可见性与脆弱性]
鲁的年表、郑的近畿军政与宋的调停席位，都是小国在制度缝隙中取得的能力。它们分别解决记录、拒压和沟通的当下难题，成本却是高度依赖王室名分、晋楚均势和贵族网络。到鲁君出奔、家臣用事，以及陈、田氏越过齐君而行动时，封邑、家臣和军政资源已在改变权力的归属：礼仪仍能命名秩序，却越来越难单独约束掌握资源的家族。
\end{turningpoint}
